% GENERATED FILE. Edit canonical lessons, then run npm run generate:books.
% canonical-source-hash: fnv1a64:f9ec31768fc663b0
% canonical-lessons: TE-C444-indhanam, TE-C444-cedipovu, TE-C444-morayincu, TE-C444-pankcaru, TE-C444-marammatu-ceyu

\chapter{Fuel and Breakdowns}
\label{ch:te-fuel-and-breakdowns}
\hlchaptermodality{444}

\begin{chapteropening}
\textbf{By the end of this chapter:} \emph{I can say fuel, to break down, to refuse to start, a puncture and to repair.}

Complete the last lesson of chapter 444: I can say fuel, to break down, to refuse to start, a puncture and to repair.
\end{chapteropening}

\section[\texorpdfstring{\te{ఇంధనం} (indhanaṁ)}{ఇంధనం (indhanaṁ)}]{\te{ఇంధనం} (indhanaṁ) — fuel}

\label{lesson:TE-C444-indhanam}



\begin{quote}
Before the new one: say the Telugu for a brake, then the Telugu for a horn.
\end{quote}

\subsection*{You'll want to know: \te{ఇంధనం}}
\textbf{\te{ఇంధనం}} — \emph{indhanaṁ} — \textquotedblleft{}fuel\textquotedblright{}.

\begin{quote}
\te{దారి} \te{మధ్యలో} \te{కారులో} \te{ఇంధనం} \te{అయిపోయింది}.
\end{quote}

\begin{quote}
\emph{dāri madhyalō kārulō indhanaṁ ayipōyindi.}
\end{quote}

\begin{quote}
Halfway there the car ran out of fuel.
\end{quote}

A word for every kind of fuel: firewood, petrol, gas. \te{అయిపోవు}, to run out, goes with it.

\begin{cousinweb}
\textbf{\te{ఇంధనం}} (\emph{indhanaṁ}) is borrowed from Sanskrit \emph{indhana}, \textquotedblleft{}lighting, kindling, fuel\textquotedblright{}, from the root \emph{indh}, \textquotedblleft{}to kindle\textquotedblright{} (CDIAL 1584).
\end{cousinweb}

\begin{grammarlens}[title={What you've built}]
One more word for getting around.
\end{grammarlens}

\subsection*{Your turn}
\begin{itemize}
\raggedright
  \item \emph{Say it:} \emph{indhanaṁ}
  \item \emph{Say it:} \emph{indhanaṁ}, once more
  \item \emph{Say it:} \emph{dāri madhyalō kārulō indhanaṁ ayipōyindi.} — Halfway there the car ran out of fuel.
  \item \emph{Recall:} say the Telugu for a brake, then the Telugu for a horn, then say \emph{indhanaṁ} again
\end{itemize}

\begin{tcolorbox}[breakable,colback=teal!4,colframe=teal!35!black,title={Before you move on}]
What does \te{ఇంధనం} mean? (\textquotedblleft{}Fuel\textquotedblright{}.) Say it once more.
\end{tcolorbox}
\section[\texorpdfstring{\te{చెడిపోవు} (ceḍipōvu)}{చెడిపోవు (ceḍipōvu)}]{\te{చెడిపోవు} (ceḍipōvu) — to break down; to go bad}

\label{lesson:TE-C444-cedipovu}



\begin{quote}
Before the new one: say the Telugu for a horn, then the Telugu for fuel.
\end{quote}

\subsection*{You'll want to know: \te{చెడిపోవు}}
\textbf{\te{చెడిపోవు}} — \emph{ceḍipōvu} — \textquotedblleft{}to break down; to go bad\textquotedblright{}.

\begin{quote}
\te{దారిలో} \te{మా} \te{బస్సు} \te{చెడిపోయింది}.
\end{quote}

\begin{quote}
\emph{dārilō mā bassu ceḍipōyindi.}
\end{quote}

\begin{quote}
Our bus broke down on the way.
\end{quote}

Past (it …): \te{చెడిపోయింది}, \emph{ceḍipōyindi}.

\te{చెడు}, bad, + \te{పోవు}, to go: gone bad. A machine breaks down this way, and food goes off this way too. Its subject is a thing, so the past shown is the third person: \te{చెడిపోయింది}, it broke down.

\begin{cousinweb}
\textbf{\te{చెడిపోవు}} (\emph{ceḍipōvu}) is native Dravidian: \te{చెడు}, \textquotedblleft{}to rot, to be spoiled\textquotedblright{} (DEDR 2760), with \te{పోవు}, \textquotedblleft{}to go\textquotedblright{} (DEDR 4572), which marks a change that has happened.
\end{cousinweb}

\begin{grammarlens}[title={What you've built}]
One more word for getting around.
\end{grammarlens}

\subsection*{Your turn}
\begin{itemize}
\raggedright
  \item \emph{Say it:} \emph{ceḍipōvu}
  \item \emph{Say it:} \emph{ceḍipōvu}, once more
  \item \emph{Say it:} \emph{dārilō mā bassu ceḍipōyindi.} — Our bus broke down on the way.
  \item \emph{Recall:} say the Telugu for a horn, then the Telugu for fuel, then say \emph{ceḍipōvu} again
\end{itemize}

\begin{tcolorbox}[breakable,colback=teal!4,colframe=teal!35!black,title={Before you move on}]
What does \te{చెడిపోవు} mean? (\textquotedblleft{}To break down; to go bad\textquotedblright{}.) Say it once more.
\end{tcolorbox}
\section[\texorpdfstring{\te{మొరాయించు} (morāyiñcu)}{మొరాయించు (morāyiñcu)}]{\te{మొరాయించు} (morāyiñcu) — to refuse to start; to balk}

\label{lesson:TE-C444-morayincu}



\begin{quote}
Before the new one: say the Telugu for fuel, then the Telugu for to break down.
\end{quote}

\subsection*{You'll want to know: \te{మొరాయించు}}
\textbf{\te{మొరాయించు}} — \emph{morāyiñcu} — \textquotedblleft{}to refuse to start; to balk\textquotedblright{}.

\begin{quote}
\te{చలికి} \te{మా} \te{కారు} \te{మొరాయించింది}.
\end{quote}

\begin{quote}
\emph{caliki mā kāru morāyiñcindi.}
\end{quote}

\begin{quote}
In the cold our car would not start.
\end{quote}

Past (it …): \te{మొరాయించింది}, \emph{morāyiñcindi}.

Said of a machine, or a stubborn animal, that will not do what it should. Its subject is a thing, so the past shown is the third person: \te{మొరాయించింది}, it would not start.

\begin{cousinweb}
\textbf{\te{మొరాయించు}} (\emph{morāyiñcu}) is a Telugu verb in -\te{ఇంచు} whose first part has an origin that is not settled.
\end{cousinweb}

\begin{grammarlens}[title={What you've built}]
One more word for getting around.
\end{grammarlens}

\subsection*{Your turn}
\begin{itemize}
\raggedright
  \item \emph{Say it:} \emph{morāyiñcu}
  \item \emph{Say it:} \emph{morāyiñcu}, once more
  \item \emph{Say it:} \emph{caliki mā kāru morāyiñcindi.} — In the cold our car would not start.
  \item \emph{Recall:} say the Telugu for fuel, then the Telugu for to break down, then say \emph{morāyiñcu} again
\end{itemize}

\begin{tcolorbox}[breakable,colback=teal!4,colframe=teal!35!black,title={Before you move on}]
What does \te{మొరాయించు} mean? (\textquotedblleft{}To refuse to start; to balk\textquotedblright{}.) Say it once more.
\end{tcolorbox}
\section[\texorpdfstring{\te{పంక్చరు} (paṅkcaru)}{పంక్చరు (paṅkcaru)}]{\te{పంక్చరు} (paṅkcaru) — a puncture}

\label{lesson:TE-C444-pankcaru}



\begin{quote}
Before the new one: say the Telugu for to break down, then the Telugu for to refuse to start.
\end{quote}

\subsection*{You'll want to know: \te{పంక్చరు}}
\textbf{\te{పంక్చరు}} — \emph{paṅkcaru} — \textquotedblleft{}a puncture\textquotedblright{}.

\begin{quote}
\te{సైకిలు} \te{టైరుకు} \te{పంక్చరు} \te{అయింది}.
\end{quote}

\begin{quote}
\emph{saikilu ṭairuku paṅkcaru ayindi.}
\end{quote}

\begin{quote}
The bicycle tyre got a puncture.
\end{quote}

\te{పంక్చరు} \te{అయింది}, it got a puncture; the shop that mends it is the \te{పంక్చరు} shop.

\begin{cousinweb}
\textbf{\te{పంక్చరు}} (\emph{paṅkcaru}) is a loan from English \emph{puncture}, with the Telugu final -\te{ఉ}.
\end{cousinweb}

\begin{grammarlens}[title={What you've built}]
One more word for getting around.
\end{grammarlens}

\subsection*{Your turn}
\begin{itemize}
\raggedright
  \item \emph{Say it:} \emph{paṅkcaru}
  \item \emph{Say it:} \emph{paṅkcaru}, once more
  \item \emph{Say it:} \emph{saikilu ṭairuku paṅkcaru ayindi.} — The bicycle tyre got a puncture.
  \item \emph{Recall:} say the Telugu for to break down, then the Telugu for to refuse to start, then say \emph{paṅkcaru} again
\end{itemize}

\begin{tcolorbox}[breakable,colback=teal!4,colframe=teal!35!black,title={Before you move on}]
What does \te{పంక్చరు} mean? (\textquotedblleft{}A puncture\textquotedblright{}.) Say it once more.
\end{tcolorbox}
\section[\texorpdfstring{\te{మరమ్మతు} \te{చేయు} (marammatu cēyu)}{మరమ్మతు చేయు (marammatu cēyu)}]{\te{మరమ్మతు} \te{చేయు} (marammatu cēyu) — to mend; to repair}

\label{lesson:TE-C444-marammatu-ceyu}



\begin{quote}
Before the new one: say the Telugu for to refuse to start, then the Telugu for a puncture.
\end{quote}

\subsection*{You'll want to know: \te{మరమ్మతు} \te{చేయు}}
\textbf{\te{మరమ్మతు} \te{చేయు}} — \emph{marammatu cēyu} — \textquotedblleft{}to mend; to repair\textquotedblright{}.

\begin{quote}
\te{మెకానిక్కు} \te{గంటలో} \te{బస్సును} \te{మరమ్మతు} \te{చేశాడు}.
\end{quote}

\begin{quote}
\emph{mekānikku gaṇṭalō bassunu marammatu cēśāḍu.}
\end{quote}

\begin{quote}
The mechanic repaired the bus within an hour.
\end{quote}

Past (I …): \te{మరమ్మతు} \te{చేశాను}, \emph{marammatu cēśānu}.

\te{మరమ్మతు}, a repair, + \te{చేయు}, to do. The thing repaired takes -\te{ని} or -\te{ను}: \te{బస్సును}.

\begin{cousinweb}
\textbf{\te{మరమ్మతు} \te{చేయు}} (\emph{marammatu cēyu}) pairs \te{మరమ్మతు}, a loan from Arabic \emph{marammat}, \textquotedblleft{}repair\textquotedblright{}, which came through Persian and Hindustani, with the native \te{చేయు}, \textquotedblleft{}to do\textquotedblright{} (DEDR 1957).
\end{cousinweb}

\begin{grammarlens}[title={What you've built}]
One more word for getting around.
\end{grammarlens}

\subsection*{Your turn}
\begin{itemize}
\raggedright
  \item \emph{Say it:} \emph{marammatu cēyu}
  \item \emph{Say it:} \emph{marammatu cēyu}, once more
  \item \emph{Say it:} \emph{mekānikku gaṇṭalō bassunu marammatu cēśāḍu.} — The mechanic repaired the bus within an hour.
  \item \emph{Recall:} say the Telugu for to refuse to start, then the Telugu for a puncture, then say \emph{marammatu cēyu} again
\end{itemize}

\begin{tcolorbox}[breakable,colback=teal!4,colframe=teal!35!black,title={Before you move on}]
What does \te{మరమ్మతు} \te{చేయు} mean? (\textquotedblleft{}To mend; to repair\textquotedblright{}.) Say it once more.
\end{tcolorbox}
